\begin{frame}{Budget Data for Comparison}
    \mytext{
    The following are illustrations showing budgets of central government, social security systems, and the city state of Bremen. 

    France’s welfare state is large relative to its economy, not unusually generous per person. Germany spends about as much per head, but from a larger GDP per head.
    }
    \mycomment{
        If you were to compare these across each other, you must account for the following:
        \begin{itemize}
            \item differences in population: Germany has roundabout 25\,\% more population than France
            \item differences in GDP per head: a richer country will have more income and can afford more expenditure per head than a poorer one.
            \item purchasing power parity: the pure numbers all are denominated in EUR, but need to be adjusted to be meaningful.
        \end{itemize}
    }
\end{frame}

\begin{frame}{Germany's State Budget (2027 draft) - Federal}
    \mycomment{
    The by far largest position in the federal budget of Germany is \enquote{Arbeit und Soziales} (BMAS), followed by defence (BMVg). The third post is \enquote{Allgemeine Finanzverwaltung}; debt service comes fourth, followed by traffic (here BMV). All other positions are smaller; the ministry for economic affairs (BMW) only receives about 1\,\% of the federal budget.

    Social security, unemployment etc.\ is only partly included here: most of it runs via contributions and not through the State budget. \enquote{Social} includes that part that does (\enquote{Grundsicherung}; subsidies to pensions, etc.).

    On the income side, about three quarters are taxes, levies etc., with another fifth \enquote{Bundesschuld}, new debt. Both sides are equal, in 2027 roundabout 555 billion Euro.
    }
    \showmytikz[1]{../shared//tikz/tikz-budgetbars}
\end{frame}

\begin{frame}{Germany's State Budget (2026) - Bremen}
    \mycomment{
    Germany is a federal state where some areas are exclusive priviledge of the Länder (the states). Most notably education, culture, and policing. So these positions do not appear on the federal level, but on the Länder level.

    According to the German constitution, tax revenue is shared between Bund, Länder and Kommunen (apart from some taxes that go directly to either Bund or Länder). As a result, the sum of all Länder budgets is be roughly the same as the federal budget. The federal state, however, borrows more, mostly outside of core budget (\enquote{Sondervermögen}). Also, substantial sums are moved between the different levels, so the budgets cannot simply be summed up.

    Bremen has the highest debt level of all German Länder by a margin, at approximately 37.000 Euro per inhabitant. (Federal debt adds another 22.000 Euro per head.) Bremen's debt is high relative to GDP, but its interest burden is moderate for now, thanks to low rates locked in during the past decade and federal consolidation aid. As old loans are refinanced at higher rates, the debt service segment will grow.
    }
    \showmytikz[2]{../shared//tikz/tikz-budgetbars}
\end{frame}

\begin{frame}{France: State Budget}
    \mycomment{For comparison, we can look at France. Of course, the numbers are not exactly comparable, since statistic methods may vary. But the picture is clear. The central government pays the teachers. Again, social insurance (pensions, health)  is mainly outside (contributions to a system independent of the state budget). However, the size of debt service and the size of new debt per year is striking.
    }
    \showmytikz[3]{../shared//tikz/tikz-budgetbars}
\end{frame}

\begin{frame}{Germany Social Insurance Budget}
    \mycomment{
    Germany's social security system is running deficits which are covered by federal funds, so this appears in the federal budget's position \enquote{social}. Unemployment scheme is among \enquote{div.}.


    }
    \showmytikz[4]{../shared//tikz/tikz-budgetbars}
\end{frame}

\begin{frame}{France Social Insurance Budget}
    \mycomment{
    In France, the social insurance system issues their own debt, and runs its own deficit. Apart from contributions, there is a general tax \enquote{contribution sociale généralisée} (CSG) on almost all income to finance the system, which could indeed be compared to a state subsidy.

    Some items still sit outside, most notably France's mandatory supplementary pensions (Agirc-Arrco, roughly 100 bio EUR) and unemployment insurance (Unédic, roughly 40 bio EUR).

    Note the central figures: the social security budget of France is larger than the state's budget. 
    }
    \showmytikz[5]{../shared//tikz/tikz-budgetbars}
    
\end{frame}